% Auto-generated by csv_to_tikz_bar.py — do not edit
\begin{figure*}[t]
\centering
\begin{tikzpicture}
\begin{axis}[
  width=\textwidth,
  height=4.0cm,
  ybar,
  bar width=6pt,
  xtick={0,...,19},
  xticklabels={
    {mg\_mload/nogr},
    {mg\_mload/by32},
    {mg\_mload/by1},
    {mg\_mload/by16},
    {snail/bench},
    {sha1\_s/5311},
    {sha1\_d/5311},
    {sha1\_s/emp},
    {sha1\_d/emp},
    {b2\_huff/8415n},
    {b2\_huff/emp},
    {b2\_shifts/8415n},
    {struct\_a/nfts},
    {swap\_m/spent},
    {swap\_m/recv},
    {swap\_m/insuf},
    {signext/zero},
    {JDEST/emp},
    {loop\_jd/emp},
    {jump\_ar/emp}
  },
  xticklabel style={rotate=45, anchor=east, font=\scriptsize\ttfamily},
  x tick label style={rotate=45, anchor=east},
  ylabel={Speedup (\%)},
  ylabel style={font=\small},
  ymin=0,
  ymax=28,
  xmin=-0.8,
  xmax=19.8,
  ytick={0,5,10,15,20,25},
  grid=major,
  grid style={dashed, gray!30},
  legend style={
    at={(0.5,1.02)},
    anchor=south,
    legend columns=4,
    font=\scriptsize,
    draw=none,
  },
  legend image code/.code={%
    \draw[#1] (0cm,-0.08cm) rectangle (0.45cm,0.08cm);%
  },
  clip=false,
]

% zero reference line
\draw[black, thin] (axis cs:-0.8,0) -- (axis cs:19.8,0);

% geometric-mean reference line (+7.24%)
\draw[blue!70!black, dashed, thick] (axis cs:-0.8,7.24) -- (axis cs:19.8,7.24);
\node[blue!70!black, font=\scriptsize, anchor=south west, fill=white, inner sep=1pt]
  at (axis cs:0,7.24) {geomean\,+7.24\%};

% group: U256-heavy
\addplot[fill=black!20, draw=black, bar shift=0pt, error bars/.cd, y dir=both, y explicit, error bar style={thin}, error mark options={rotate=90, mark size=1.5pt}]
  table[x=x, y=y, y error minus=ym, y error plus=yp] {
    x y ym yp
    0 7.2549 0.8283 0.8210
    1 6.9940 1.0760 1.0637
    2 6.9362 0.7959 0.7892
    3 6.0337 0.7583 0.7522
    4 2.6203 0.7584 0.7526
  };
\addlegendentry{U256-heavy}

% group: Hash
\addplot[pattern=north east lines, draw=black, bar shift=0pt, error bars/.cd, y dir=both, y explicit, error bar style={thin}, error mark options={rotate=90, mark size=1.5pt}]
  table[x=x, y=y, y error minus=ym, y error plus=yp] {
    x y ym yp
    5 24.8066 0.6719 0.6660
    6 23.5105 0.7462 0.7390
    7 20.2360 1.0368 1.0235
    8 19.5875 0.7839 0.7763
    9 11.5442 1.0606 1.0481
    10 10.4090 0.7829 0.7761
    11 4.4577 0.8946 0.8863
  };
\addlegendentry{Hash}

% group: Algorithmic
\addplot[pattern=crosshatch, draw=black, bar shift=0pt, error bars/.cd, y dir=both, y explicit, error bar style={thin}, error mark options={rotate=90, mark size=1.5pt}]
  table[x=x, y=y, y error minus=ym, y error plus=yp] {
    x y ym yp
    12 17.0420 0.7116 0.7056
    13 4.7468 0.9761 0.9662
    14 3.8131 0.9158 0.9072
    15 2.4458 0.9211 0.9124
    16 0.6157 1.3616 1.3432
  };
\addlegendentry{Arithmetic-Logic}

% group: Control-flow
\addplot[pattern=dots, draw=black, bar shift=0pt, error bars/.cd, y dir=both, y explicit, error bar style={thin}, error mark options={rotate=90, mark size=1.5pt}]
  table[x=x, y=y, y error minus=ym, y error plus=yp] {
    x y ym yp
    17 25.7880 0.6690 0.6630
    18 4.8763 0.6924 0.6873
    19 1.5885 1.2836 1.2671
  };
\addlegendentry{Control-flow}

\end{axis}
\end{tikzpicture}
\par\vspace{-4pt}
{\scriptsize \emph{Labels:} mg=\texttt{memory\_grow}, b2=\texttt{blake2b}, sha1\_s=\texttt{sha1\_shifts}, sha1\_d=\texttt{sha1\_divs}, struct\_a=\texttt{struct\_array}, swap\_m=\texttt{swap\_method}, signext=\texttt{signextend}, JDEST=\texttt{JUMPDEST\_n0}, loop\_jd=\texttt{loop\_jumpdest}, jump\_ar=\texttt{jump\_around}, snail=\texttt{snailtracer}.}
\vspace{-4pt}
\caption{Positive-speedup benchmark configurations from enabling the
peephole rewrite rule set over the upstream-main baseline (error bars
are 95\% bootstrap confidence intervals, $B{=}10\,000$). The horizontal
line at zero marks the no-change baseline; the dashed line at
\textbf{+7.24\%} marks the 27-benchmark geometric-mean speedup. From
left to right, the benchmarks are grouped into four workload clusters:
U256-heavy, Hash, Arithmetic-Logic, and Control-flow.}
\label{fig:speedup}
\end{figure*}
